\documentclass[11pt]{article}

\usepackage[T1]{fontenc}
\usepackage{amsmath,amssymb,amsthm}
\usepackage[margin=1in]{geometry}
\usepackage[hidelinks]{hyperref}

\newtheorem{theorem}{Theorem}
\newtheorem{proposition}[theorem]{Proposition}
\newtheorem{lemma}[theorem]{Lemma}
\newtheorem{corollary}[theorem]{Corollary}
\theoremstyle{definition}
\newtheorem{definition}[theorem]{Definition}
\newtheorem{remark}[theorem]{Remark}

\newcommand{\Z}{\mathbb Z}
\newcommand{\Q}{\mathbb Q}
\newcommand{\R}{\mathbb R}
\newcommand{\mass}{\operatorname{m}}

\title{The Rank-One Mass-Three Radius for Five-Letter Additive Models Is 33}
\author{Xiaojian Zeng\\
  \small PBC School of Finance, Tsinghua University, Beijing 100190, China\\
  \small \texttt{xiaojianzeng@foxmail.com}}
\date{September 2026}

\begin{document}
\maketitle

\begin{abstract}
For an increasing real five-letter alphabet, let $M_3$ be the least
diameter of an increasing integer alphabet preserving exactly all equalities
between two equal-length sums of at most three letters.  We prove that
$M_3\leq 33$ whenever the rational span of the short relations has rank one
and its primitive generator has mass three.  The bound is sharp.  Moreover,
up to changing the sign of the generator, equality occurs only for the
relation directions
\[
  (1,-3,0,0,2)\quad\text{and}\quad(2,0,0,-3,1).
\]
The proof reduces the real problem to an exact finite classification.  There
are 160 primitive mass-three directions modulo sign, 65 of which meet the
strict order cone.  A complete census of the $\binom{33}{4}=40{,}920$
normalized integer alphabets of diameter at most 33 supplies an exact-profile
model for every feasible direction.  A second implementation regenerates
the relation domain and repeats the census.  A separate 165-core census
establishes the sharp lower bound.
\end{abstract}

\section{Introduction}

An order-preserving Freiman model retains the order of the letters and the
specified additive equalities.  The order-two formulation was studied by
Amirkhanyan, Bush and Croot~\cite{AmirkhanyanBushCroot2016}; the equality-only
higher-order invariant used here is the zero-profile formulation of
Zhang~\cite{Zhang2026Sharp}.  In contrast with a sign-preserving addition
table, signs of unequal sums are allowed to change.

Zhang's mass-sensitive rank-one estimate gives diameter at most
$\mu q(q+1)$ for a primitive relation of mass $\mu$~\cite[Proposition
6.5]{Zhang2026Sharp}.  At $(q,\mu)=(3,3)$ this is 36.  On the other hand, the
cluster profile in~\cite[Corollary 9.2]{Zhang2026Sharp} has exact diameter
$3L_4(3)=33$.  The value $L_4(3)=11$ is the $q=3$ instance of Shin's exact
four-mark relation-free formula~\cite[Theorem 10.3]{Shin2026Spectra}.
Zhang's Question 12.3 asks for the rank-one maxima by primitive
mass~\cite[Question 12.3]{Zhang2026Sharp}.  The following theorem determines
the case $(q,\mu)=(3,3)$.

\begin{theorem}\label{thm:main}
Let $T=(t_0<t_1<t_2<t_3<t_4)$ be a real alphabet.  Suppose that the rational
span of its balanced mass-at-most-three zero relations has rank one and that
the primitive generator of this span has mass three.  Then
\[
                         M_3(T)\leq 33.
\]
The maximum over all such alphabets is exactly 33.  Up to replacing the
primitive generator by its negative, the only relation directions attaining
33 are
\[
  \delta_{\mathrm c}=(1,-3,0,0,2)
  \quad\text{and}\quad
  \delta_{\mathrm c}^{\mathrm{rev}}=(2,0,0,-3,1).
\]
\end{theorem}

The comparison above is with the cited versions inspected in September
2026.  Because the originating preprint appeared immediately before this
work, the theorem is not accompanied by a claim of global priority.

\section{Profiles and the order cone}

For $a=(a_0,\ldots,a_4)\in\Z^5$ with $\sum_i a_i=0$, define its mass by
\[
 \mass(a)=\sum_i\max(a_i,0)
          =\sum_i\max(-a_i,0)=\frac12\sum_i|a_i|.
\]
The mass-at-most-three zero profile of $T$ is
\[
 Z_3(T)=\left\{a\in\Z^5:\ \sum_i a_i=0,\ \mass(a)\leq3,
                         \ a\cdot T=0\right\}.
\]
An increasing integer alphabet $S=(s_0<\cdots<s_4)$ is a $3$-model of $T$
when $Z_3(S)=Z_3(T)$.  Its diameter is $s_4-s_0$, and $M_3(T)$ is the least
such diameter.  Thus equality must be preserved in both directions; no
condition is imposed on the signs of nonzero sum differences.

\begin{lemma}[Rank-one reduction]\label{lem:profile}
Under the hypotheses of Theorem~\ref{thm:main}, if $\delta$ is the primitive
balanced generator of the relation line, then
\[
                         Z_3(T)=\{0,\delta,-\delta\}.
\]
\end{lemma}

\begin{proof}
Primitivity gives $\Q\delta\cap\Z^5=\Z\delta$.  Since the relation line is
spanned by $Z_3(T)$, it contains a nonzero member $k\delta$ of the profile,
where $k\in\Z\setminus\{0\}$.  Homogeneity of mass and
$\mass(\delta)=3$ force $|k|=1$.  The same argument applies to every nonzero
member of the profile.
\end{proof}

For a balanced row $a$, put
\[
                         A_j(a)=\sum_{i=j}^4a_i
                         \qquad(1\leq j\leq4).
\]

\begin{lemma}[Tail-sum criterion]\label{lem:tails}
For a nonzero balanced row $a$, the hyperplane $a\cdot T=0$ contains a strictly increasing real alphabet if
and only if the four tail sums $A_1(a),\ldots,A_4(a)$ contain both a positive
and a negative value.
\end{lemma}

\begin{proof}
Set $g_j=t_j-t_{j-1}$.  Balance gives the exact summation-by-parts identity
\[
                         a\cdot T=\sum_{j=1}^4 A_j(a)g_j.
\]
If all $g_j>0$ and the right side is zero, the nonzero tail row must have both
signs.  Conversely, let
\[
 P=\sum_{A_j>0}A_j,
 \qquad N=-\sum_{A_j<0}A_j.
\]
Both are positive.  Choose $g_j=N$ when $A_j>0$, $g_j=P$ when $A_j<0$, and
$g_j=1$ when $A_j=0$.  Then every gap is positive and the displayed sum is
$PN-NP=0$.
\end{proof}

The criterion concerns intersection with the strict order cone.  It also
gives realizability as a rank-one short profile: inside the resulting
relative open set, avoid the finitely many hyperplanes belonging to short
rows not proportional to $a$.  None of them contains the whole set.  When
$a$ is primitive of mass three, the only proportional integral short rows
are $\pm a$.

\section{The exact finite classification}

Fix the sign of every nonzero integer row by requiring its first nonzero
coordinate to be positive.  For a feasible primitive mass-three direction
$\delta$, let $D(\delta)$ be the least diameter of an increasing integer
alphabet whose canonical nonzero profile is the singleton $\{\delta\}$.

\begin{proposition}[Finite classification]\label{prop:census}
There are 160 primitive balanced mass-three directions modulo sign.  Exactly
65 satisfy the tail-sum criterion.  Every one of these 65 directions has
$D(\delta)\leq33$, and
\[
 \max_{\delta}D(\delta)=33.
\]
The directions attaining this maximum are exactly
$\delta_{\mathrm c}$ and $\delta_{\mathrm c}^{\mathrm{rev}}$.
\end{proposition}

\begin{proof}[Exact certificate proof]
Every coordinate of a balanced row of mass at most three lies in
$[-3,3]$.  Hence the coefficient box $[-3,3]^5$, together with balance and
the mass test, is a complete domain.  For reference, the number of oriented
rows of exact mass $p$ is
\[
 N_p=\sum_{\substack{r,s\geq1\\r+s\leq5}}
   \binom5r\binom{5-r}s
   \binom{p-1}{r-1}\binom{p-1}{s-1}.
\]
This gives $(N_1,N_2,N_3)=(20,110,340)$, so there are
$(20+110+340)/2=235$ nonzero short rows modulo sign.  The only nonprimitive
mass-three rows have the form $3(e_i-e_j)$; removing their 20 orientations
leaves $(340-20)/2=160$ primitive directions modulo sign.

For each of these 160 rows, the four integer tail sums are evaluated
literally.  Sixty-five have both signs.  For the model census, translation
reduces every integer alphabet of diameter $D$ to
\[
                 S=(0<s_1<s_2<s_3<D).
\]
The complete domain through diameter 33 therefore has size
\[
 \sum_{D=4}^{33}\binom{D-1}{3}=\binom{33}{4}=40{,}920.
\]
For each tuple, all 235 canonical short rows are evaluated with integer
arithmetic.  The tuple is assigned to $\delta$ precisely when its entire
canonical nonzero zero profile is $\{\delta\}$.  Diameters are traversed in
increasing order, so the first assigned tuple is a minimum model.  Every one
of the 65 feasible directions is assigned; the largest first diameter is 33,
and the complete list at that diameter consists of the two stated rows.

The accompanying producer performs exactly this coefficient-box census.  A
separate checker does not import its generators: it reconstructs all short
rows as differences of unordered multisets of one, two, or three labels, and
it reconstructs mass-three directions from pairs of disjoint weak
compositions of 3.  It then repeats all 40,920 model tests and checks every
stored minimum and occurrence count.  Both implementations use only exact
integers.  This separation of the mathematical reduction, stored output,
and independent exact verification is also used in other finite
combinatorial certificate work; see, for example,
Zeng, Liu, Ratnavelu and Ong~\cite{ZengEtAl2026FiniteCertificates}.
\end{proof}

\begin{corollary}\label{cor:upper}
Every alphabet in Theorem~\ref{thm:main} has a $3$-model of diameter at most
33.  Moreover, if its primitive direction is not one of the two directions
in Theorem~\ref{thm:main}, its minimum diameter is at most 32.
\end{corollary}

\begin{proof}
By Lemma~\ref{lem:profile}, the profile is determined by its primitive
direction.  Lemma~\ref{lem:tails} places that direction among the 65 feasible
rows.  Proposition~\ref{prop:census} supplies an integer alphabet with
exactly the same profile and gives the stated minima.
\end{proof}

\section{Sharpness from a four-letter core}

The complete census already records the minimum of the cluster direction.
For clarity, we isolate a smaller certificate proving its lower bound.

\begin{lemma}[Cluster reduction]\label{lem:cluster}
An increasing integer alphabet beginning at zero has canonical profile
$\{\delta_{\mathrm c}\}$ if and only if it has the form
\[
 S=(0,2c,2c+a,2c+b,3c),\qquad 0<a<b<c,
\]
and the four-letter core $U=(0,a,b,c)$ has no nonzero balanced relation of
mass at most three.
\end{lemma}

\begin{proof}
The relation $\delta_{\mathrm c}\cdot S=0$ is
$-3s_1+2s_4=0$.  Since $\gcd(2,3)=1$, this forces $s_1=2c$ and $s_4=3c$;
strict order gives the displayed form.  Equal-length relations supported on
the last four letters are unchanged when $2c$ is subtracted from every
letter, so they are exactly the relations on $U$.

It remains to examine a reduced relation that uses the initial zero.  Orient
it so that $z\geq1$ zeros occur on one side, and let its mass be $p\leq3$.
After cancellation the other side has $p$ nonzero terms, each at least
$2c$, while the zero side has only $p-z$ nonzero terms, each at most $3c$.
Equality would require
\[
                         3c(p-z)\geq2cp,
\]
or $p\geq3z$.  Hence $p=3$ and $z=1$.  Equality throughout forces the two
remaining terms on the zero side to be $3c$ and all three terms on the other
side to be $2c$.  This is precisely $0+2s_4=3s_1$, the cluster relation.
Thus no additional relation involving zero is possible.
\end{proof}

Because $U$ begins at zero, it is relation-free through mass three exactly
when its 20 unordered three-term sums are pairwise distinct: any shorter
equality can be padded with zeros, and a collision gives a short balanced
relation after cancellation.

\begin{proposition}[Core census]\label{prop:core}
No integer core $(0,a,b,c)$ with $0<a<b<c\leq10$ is relation-free through
mass three.  At $c=11$, the complete list of relation-free cores is
\[
 (0,1,7,11),\quad(0,1,8,11),\quad
 (0,3,10,11),\quad(0,4,10,11).
\]
\end{proposition}

\begin{proof}[Exact certificate proof]
For each $3\leq c\leq11$, enumerate the $\binom{c-1}{2}$ choices
$1\leq a<b<c$ and test the 20 integer triple sums for a collision.  The
domain contains
\[
                    \sum_{c=3}^{11}\binom{c-1}{2}
                    =\binom{11}{3}=165
\]
cores.  Direct enumeration gives zero successes through $c=10$ and the four
displayed successes at $c=11$.  The independent checker regenerates the
four-coordinate relation rows and repeats the test.
\end{proof}

For $\delta_{\mathrm c}$, Lemma~\ref{lem:cluster} and
Proposition~\ref{prop:core} imply $c\geq11$, so every exact-profile model has
diameter $3c\geq33$.  The first core in Proposition~\ref{prop:core} yields
\[
                         (0,22,23,29,33),
\]
which has exactly the cluster profile.  A real representative is
$(0,2,2+\theta,2+\eta,3)$, where $0<\theta<\eta<1$ and
$1,\theta,\eta$ are linearly independent over $\Q$; its only nonzero short
relations are the two signs of $\delta_{\mathrm c}$.  Reflection gives the
reversal direction.  Together with Corollary~\ref{cor:upper} and the complete
extremal list in Proposition~\ref{prop:census}, this proves all assertions of
Theorem~\ref{thm:main}.

\section*{Data, code, and assistance disclosure}

The ancillary evidence consists of the producer
\texttt{enumerate\_mass3.py}, its JSON classification, the independently
written checker \texttt{check\_mass3\_certificate.py}, and the checker's JSON
audit.  The programs use the Python standard library and exact integer
arithmetic; no floating-point comparison, random seed, or numerical tolerance
enters a theorem-bearing calculation.  Their finite domains and acceptance
predicates are stated in the proof above.

OpenAI Codex assisted with the bounded literature search, finite-search
design and implementation, organization of a separate mathematical audit,
and initial manuscript preparation.  Language-model output is not used as
mathematical evidence: the result rests on the displayed reductions and the
replayable exact enumerations.  No proof assistant was used.

\bibliographystyle{abbrv}
\bibliography{references}

\end{document}
